%Abstract Page

\hbox{\ } \vspace{.7in}
\renewcommand{\baselinestretch}{1}
\small \normalsize

\begin{center}
\large{{ABSTRACT}}

\vspace{3em}

\end{center}
\hspace{-.15in}
\begin{tabular}{ll}
Title of Dissertation:    & {\large ENGINEERING EXCITON INTERACTIONS}\\
&                     {\large  IN VAN DER WAALS HETEROSTRUCTURE} \\
\ \\
&                          {\large  Liuxin Gu} \\
&                           {\large Doctor of Philosophy, 2025} \\
\ \\
Dissertation Directed by: & {\large  Professor You Zhou} \\
&               {\large  Department of Materials Science and Engineering } \\
\end{tabular}

\vspace{3em}

\renewcommand{\baselinestretch}{2}
\large \normalsize

Light-matter interaction is fundamental to both classical and quantum nonlinear optics. 
In general, photons do not interact with each other except for extreme conditions (strong spatial confinement). 
With strong nonlinear interaction, photon blockade can occur, wherein a photon coupled to the system inhibits the absorption of subsequent photons, analogous to Coulomb blockade. 
Realizing such strong interactions typically requires a high photon density and an interaction strength determined by the material properties.

2D materials with atomic thickness exhibit exceptionally strong light–exciton interactions. 
In these systems, exciton–exciton interactions give rise to nonlinear optical responses, observable as a blue shift in the exciton resonance with increasing pulsed laser excitation. 
The interaction strength, $g_{exc}$, defined as the exciton energy shift as exciton density, is dominated by exchange interactions on the order of $\sim$ $E_B \cdot\ R_b^2$, where $E_B$ is the exciton binding energy and $R_b$ is the exciton Bohr radius. 
A small Bohr radius usually leads to a tighter binding of the exciton. 
$g_{exc}$ is usually on the order of $\sim \mu eV \cdot\ \mu m^2$, whether in TMDs or GaAs quantum wells. 

In this thesis, we demonstrated that introducing free carriers into 2D materials can significantly enhance these interactions. 
In trilayer WSe$_2$, excitons can bind with carriers to form Fermi polaron, 
which exhibit large optical nonlinear behavior due to the interaction between exciton and free carriers. 
Remarkably, this interaction strength can reach  $\sim 2 meV \cdot\ \mu m^2$, nearly 1000 times stronger than exciton-exciton interactions. 

% In addition, by coupling these Fermi polarons to optical 1D microcavity,
% we realized polaron-polariton with distinct nonlinear behavior as the exciton polariton in the trilayer system. 
% Though the number of photons needed is as the same amount for the nonlinear shift as in the plain trilayer WSe$_2$, it can be further decreased with some engineering on enhancing the cavity quality.
% With higher quality of cavity, the formed polaritons are expected to exhibit longer lifetimes and even stronger nonlinear interactions, potentially making polariton blockade achievable in this system.
By coupling these Fermi polarons to a one-dimensional optical microcavity, we realized polaron--polaritons with nonlinear behavior distinct from exciton--polaritons in the same trilayer system. 
The photon number needed to induce a given nonlinear shift is comparable to that in bare trilayer WSe$_2$, but can be reduced by engineering a higher cavity quality factor $Q$ and smaller mode volume $V$. 
With improved cavity quality, the expected longer polariton lifetimes and enhanced per-polariton nonlinear shift could push the polariton to the blockade regime achievable in this platform.

Another route towards the photon blockade is by exciton confinement. Confining excitons below subwavelength scales discretizes their spectrum and strengthens exciton--exciton interactions. 
Howeverm unlike electrons readily controlled by gate fields,
imposing strong nanoscale potentials on excitons to enable quantum confinement has proven challenging.  
Here, we realize such confinement by imprinting the electrostatic superlattice of twisted hBN onto an adjacent monolayer MoSe$_2$.
By correlating real-space PFM map with optical spectra, we resolve excitons trapped in a nanoscale, one-dimensional potential formed by strong in-plane fields at moiré domain boundaries.



